$\Rightarrow p \mid r$
$\Rightarrow r \geq p$
$\Rightarrow |Z(G)| \geq p$

- Theorem: A group of order $p^2$, $p$ is prime, is abelian.

Proof: Let $|G| = p^2$
$\Rightarrow |Z(G)| \geq p$
So, $|Z(G)| = p$ or $p^2$
If $|Z(G)| = p$ then $\frac{|G|}{|Z(G)|} = p$

$\Rightarrow G/Z(G)$ is cyclic
$\Rightarrow G$ is abelian
$\Rightarrow G = Z(G) \rightarrow |G| = |Z(G)|$
$\therefore |Z(G)| \neq p$ (contradiction)
So $|Z(G)| = p^2 = |G|$
$\Rightarrow G = Z(G)$
$\Rightarrow G$ is abelian.

Rem. If $H \subseteq G$ s.t. $H \subseteq Z(G)$ then $H \triangleleft G$.

- Theorem: Let $\sigma$ be a permutation in $S_n$ and let $m_1, m_2, \dots, m_r$ be the distinct integers which appear in the cycle type of $\sigma$ (including $1$-cycles). For each $i \in \{1, 2, \dots, s\}$ assume $\sigma$ has $k_i$ cycles of length $m_i$ so that $\sum_{i=1}^s k_i m_i = n$. Then number of conjugate of $\sigma$ is $\frac{n!}{(k_1! m_1^{k_1})(k_2! m_2^{k_2}) \dots (k_s! m_s^{k_s})}$
